\chapter{基础知识}

本章建立理解后续设计所需的理论基础。内容从格与模运算开始，依次说明格困难问题、环对象的
系数与频域表示、离散Gaussian、格陷门及低秩线性代数，再介绍HIBE、IBS和IB-ME的基本语法。
这些概念既用于描述密码方案，也用于解释后文对象尺寸、采样成本和结构化计算的来源。

\section{数学符号与格基础}

\subsection{向量、矩阵与模运算}

记整数集合为$\mathbb Z$，模$q$整数环为$\mathbb Z_q=\mathbb Z/q\mathbb Z$。向量使用粗体
小写字母表示，矩阵使用大写字母表示。对$x\in\mathbb Z_q$，程序既可以保存区间$[0,q)$中的
非负代表元，也可以在需要计算范数时转换为中心区间$(-q/2,q/2]$中的代表元；两种表示必须在
接口处明确区分。

对向量$\mathbf x=(x_1,\ldots,x_n)$，报告主要使用欧氏范数和无穷范数
\[
 \lVert\mathbf x\rVert_2=\sqrt{\sum_{i=1}^{n}x_i^2},
 \qquad
 \lVert\mathbf x\rVert_\infty=\max_i|x_i|.
\]
签名方案通常要求响应向量满足给定范数界；加密方案则要求误差足够小，使解密时能够区分编码
区间。各方案以公开参数$pp$记录维数、模数、范数界及采样配置。

表~\ref{tab:foundation-symbols}列出后文频繁使用的符号。

\begin{table}[H]
\centering\small
\caption{基础数学符号}
\label{tab:foundation-symbols}
\begin{tabularx}{\textwidth}{@{}p{.22\textwidth}X@{}}
\toprule
\rowcolor{black!8} 符号 & 含义 \\
\midrule
$n,m,d,k,\ell$ & 格维数、矩阵宽度、环次数或模块秩等整数参数，具体含义由方案给出 \\
$q$ & 模数；$\mathbb Z_q$表示模$q$整数环 \\
$A,\mathbf a$ & 公开矩阵或公开向量 \\
$\mathbf s,\mathbf e$ & 秘密向量与小误差向量 \\
$\lVert\cdot\rVert_2,\lVert\cdot\rVert_\infty$ & 欧氏范数与无穷范数 \\
$R_q$ & 多项式商环$\mathbb Z_q[x]/(f(x))$ \\
$D_{\Lambda,s,\mathbf c}$ & 格$\Lambda$上中心为$\mathbf c$、宽度为$s$的离散Gaussian分布 \\
\bottomrule
\end{tabularx}
\end{table}

\subsection{格、$q$元核格与目标陪集}

给定线性无关向量组成的基$B=[\mathbf b_1,\ldots,\mathbf b_n]\in\mathbb R^{m\times n}$，
由$B$生成的格为
\[
  \Lambda(B)=\{B\mathbf z:\mathbf z\in\mathbb Z^n\}.
\]
同一个格可以具有许多不同的基，基的几何质量会影响求近向量和Gaussian采样等算法。早期格陷门通常采用一组短而较好正交的格向量；后来的结构化陷门采用短矩阵等秘密辅助信息。公开对象通常表示为模矩阵或环向量
\cite{gpv2008,micciancio2012trapdoors}。

模$q$矩阵$A\in\mathbb Z_q^{n\times m}$定义从$\mathbb Z^m$到$\mathbb Z_q^n$的映射
$\mathbf x\mapsto A\mathbf x\bmod q$。该映射的核为$q$元核格
\[
 \Lambda_q^\perp(A)=\{\mathbf x\in\mathbb Z^m:A\mathbf x=\mathbf0\pmod q\},
\]
给定目标$\mathbf u$时，相应解集为
\[
 \Lambda_q^{\mathbf u}(A)=\{\mathbf x\in\mathbb Z^m:A\mathbf x=\mathbf u\pmod q\}.
\]
若$\mathbf x_0$是其中一个解，则
\[
 \Lambda_q^{\mathbf u}(A)=\mathbf x_0+\Lambda_q^\perp(A).
\]
因此，$\Lambda_q^{\mathbf u}(A)$是核格的一个目标陪集，由核格平移$\mathbf x_0$得到。身份基发钥和格签名经常需要在该陪集中找到满足分布与范数要求
的短向量。陷门帮助合法持有者高效完成满足目标方程的短原像采样。

\section{格上困难问题及结构变体}

\subsection{LWE问题}

学习带误差问题（Learning With Errors，LWE）由Regev系统提出并用于构造公钥密码
\cite{regev2005lwe}。取均匀矩阵$A\in\mathbb Z_q^{m\times n}$、秘密$\mathbf s\in\mathbb
Z_q^n$和小误差$\mathbf e\in\mathbb Z^m$，观察值为
\[
  \mathbf b=A\mathbf s+\mathbf e\pmod q.
\]
搜索型LWE要求从$(A,\mathbf b)$恢复$\mathbf s$；判定型LWE要求区分$(A,\mathbf b)$与均匀
样本。若没有误差，线性方程可以直接求解；小误差使恢复秘密成为格上的困难问题。格加密通常将
明文或封装种子叠加在这类带误差关系上，合法密钥可以消去主线性项并解码小扰动。

\subsection{SIS问题}

短整数解问题（Short Integer Solution，SIS）与Ajtai建立的平均情形格问题有关
\cite{ajtai1996}。给定均匀矩阵$A\in\mathbb Z_q^{n\times m}$，目标是找到非零短向量
$\mathbf z$，满足
\[
 A\mathbf z=\mathbf0\pmod q,
 \qquad 0<\lVert\mathbf z\rVert_2\leq\beta.
\]
非齐次版本给出目标$\mathbf u$，要求寻找短向量$\mathbf z$使$A\mathbf z=\mathbf u\pmod q$。
格签名常把身份和消息哈希为目标向量，再由陷门采样满足方程的短原像。验证者检查模方程和范数，
而不知道陷门的参与者难以为新目标构造合格短向量。

\subsection{环格、模格与NTRU结构}

普通矩阵格直接保存$A$及其向量，表达灵活，但公开矩阵和密钥可能较大。环格把对象组织为商环
\[
 R_q=\mathbb Z_q[x]/(x^N+1)
\]
中的多项式，后文使用的参数通常取$N$为2的幂。每个环元素都可唯一写成
$a_0+a_1x+\cdots+a_{N-1}x^{N-1}$，因而只需保存$N$个模$q$系数。由于商环中$x^N=-1$，
乘积超过$N-1$次的项折回时改变符号；若$c(x)=a(x)b(x)$，则
\[
 c_k=\sum_{i+j=k}a_i b_j-\sum_{i+j=k+N}a_i b_j\pmod q,
 \qquad 0\leq k<N.
\]
这就是负循环乘法。一个$r\times t$环矩阵包含$Nrt$个系数，其作用等价于一个更大的结构化
系数矩阵，但实际计算没有必要把它展开。Ring-LWE和Ring-SIS由此获得紧凑表示和快速多项式
运算\cite{lpr2013ring}。

模格在$R_q$上使用低秩向量或矩阵，在普通矩阵格的灵活性和环格的结构性之间折中；对应问题通常
称为Module-LWE与Module-SIS\cite{langlois2015module}。ML-DSA和本作品采用的多种原生IBS
属于这一技术路线。

NTRU结构同样在多项式环上运算，通过短多项式关系形成结构化格。该结构可以用于加密、签名和
陷门委派。LATTE利用NTRU结构完成HIBE，Falcon利用NTRU格完成普通数字签名。这些格结构分别对应不同的对象组织方式与困难问题来源。

\subsection{环伴随与频域表示}

对实系数环元素$f(x)=\sum_{i=0}^{N-1}f_i x^i$，定义环上的伴随
\[
 f^\star(x)=f(x^{-1})\bmod(x^N+1)
 =f_0-\sum_{i=1}^{N-1}f_i x^{N-i}.
\]
对环矩阵，$T^\star$表示先转置再对每个元素取上述伴随。取
$\omega_j=\exp((2j+1)\pi\mathrm i/N)$，在$x^N+1$的复根上计算$f(\omega_j)$，即可得到负循环
Fourier表示。此时环乘法在每个频率上变为复数乘法，并且
$f^\star(\omega_j)=\overline{f(\omega_j)}$；环矩阵的伴随相应变为复矩阵的共轭转置。实系数对象的
频率值成共轭对，因此实现可以只保存一半频率并由对称性恢复其余部分。

系数表示适合精确模运算与对象编码，频域表示适合卷积、范数估计和小矩阵分解。实现通过逆变换和离散化恢复整数对象，再执行精确模方程检查。

\begin{table}[H]
\centering\small
\caption{本作品涉及的格结构及其一般工程特征}
\label{tab:lattice-families}
\begin{tabularx}{\textwidth}{@{}p{.18\textwidth}p{.25\textwidth}X@{}}
\toprule
\rowcolor{black!8} 结构 & 典型对象 & 一般工程特征 \\
\midrule
普通矩阵格 & $\mathbb Z_q$上的矩阵与整数向量 & 表达直接，目标参数下矩阵和陷门对象可能较大 \\
环格 & 商环多项式与环向量 & 可利用卷积和频域结构，需处理环采样与数值变换 \\
模格 & 小秩环向量或环矩阵 & 在结构复用与方案灵活性之间折中 \\
NTRU结构格 & 满足短多项式关系的结构化基 & 对象较紧凑，可支持快速签名、加密和陷门委派 \\
\bottomrule
\end{tabularx}
\end{table}

\section{离散Gaussian与格陷门}

\subsection{离散Gaussian采样}

对格$\Lambda$、中心$\mathbf c$和宽度$s>0$，连续Gaussian权重写为
\[
 \rho_{s,\mathbf c}(\mathbf x)=
 \exp\!\left(-\pi\frac{\lVert\mathbf x-\mathbf c\rVert_2^2}{s^2}\right).
\]
令$\rho_{s,\mathbf c}(\Lambda)=\sum_{\mathbf y\in\Lambda}\rho_{s,\mathbf c}(\mathbf y)$，离散
Gaussian分布的完整概率为
\[
 \Pr[X=\mathbf x]
 =\frac{\rho_{s,\mathbf c}(\mathbf x)}{\rho_{s,\mathbf c}(\Lambda)},
 \qquad \mathbf x\in\Lambda.
\]
在GPV与MP12采用的$\exp(-\pi\lVert\cdot\rVert^2/s^2)$约定下，相应一维连续Gaussian的标准差为$s/\sqrt{2\pi}$。若另一资料使用
$\exp(-\lVert\cdot\rVert^2/(2\sigma^2))$，则两种参数满足$s=\sqrt{2\pi}\sigma$，使用时据此换算\cite{gpv2008,micciancio2012trapdoors}。

当采样空间是$q$元核格的目标陪集时，输出还满足给定模方程。宽度过小可能不满足方案需要的
平滑或统计条件，宽度过大又会增大签名范数或解密误差，因此必须采用具体方案给出的参数和采样
方法。模方程与范数检查分别核对样本的代数关系和大小，采样算法则规定输出分布。

实际实现中，普通矩阵格、环格和NTRU结构使用的采样器并不相同。采样过程可能包含扰动采样、
Gram--Schmidt信息、FFT／NTT变换、Babai反演或逐系数拒绝采样。采样器按目标离散Gaussian分布生成样本，并对输出执行精确模方程与范数检查。

\subsection{陷门生成、委派与短原像采样}

格陷门是一份能够帮助合法持有者完成短原像采样的秘密辅助信息。典型算法可以概括为：

\begin{enumerate}
  \item $\mathsf{TrapGen}$生成统计上接近均匀的公开矩阵$A$及短陷门$T_A$；
  \item $\mathsf{SamplePre}(A,T_A,\mathbf u)$为目标$\mathbf u$采样短向量$\mathbf x$，使
        $A\mathbf x=\mathbf u\pmod q$；
  \item 委派算法把已有陷门扩展到包含新身份块的公开矩阵，同时保持可用的短基或结构化陷门。
\end{enumerate}

各方案分别采用具体算法实现上述三项功能。GPV采样器以短格基或短满秩向量集为
秘密几何信息；Micciancio--Peikert（MP12）直接使用一份小范数矩阵作为结构化陷门\cite{gpv2008,micciancio2012trapdoors}。Lai等人又将MP12型陷门扩展到
理想格，以环元素代替相应的结构化矩阵块\cite{lai2015trapdoors}。

\subsection{gadget矩阵与结构化陷门}

MP12框架从一个容易求逆和采样的公开gadget矩阵出发。令$k=\lceil\log_2q\rceil$、
$\mathbf g^T=(1,2,\ldots,2^{k-1})$，则常见二进制gadget为
\[
 G=I_n\otimes\mathbf g^T\in\mathbb Z_q^{n\times nk}.
\]
$G$把$n$个低维分解问题并排组织起来，可以高效完成位分解、求逆和目标陪集预像采样。采用其他
进制或环结构时，gadget的具体形式会改变，但“公开、规则、便于求预像”的作用保持不变。

对于公开矩阵$A$，MP12型$G$-陷门是一份短矩阵$R$，使
\[
 A\begin{bmatrix}R\\I\end{bmatrix}=HG\pmod q,
\]
其中$H$是可逆标签。生成时可取
\[
 A=[\bar A\mid HG-\bar A R].
\]
合法持有者先把一般目标转换为gadget上的低维预像问题，再结合扰动采样得到$A\mathbf x=\mathbf u
\pmod q$的短解。HIBE的委派进一步为扩展后的身份矩阵生成新陷门；原生格IBS则用短原像作为
身份密钥或签名。图~\ref{fig:preimage-sampling-foundation}归纳各对象的关系。

\begin{figure}[H]
\centering
\begin{tikzpicture}[node distance=7mm and 8mm]
 \node[flow,text width=37mm] (relation) {公开关系\\$A\mathbf x=\mathbf u\bmod q$};
 \node[flow,text width=37mm,right=of relation] (coset) {目标陪集\\$\mathbf x_0+\Lambda_q^\perp(A)$};
 \node[check,text width=38mm,below=of relation] (trapdoor) {秘密陷门与公开gadget\\提供高效采样路径};
 \node[check,text width=38mm,right=of trapdoor] (parameter) {Gaussian参数\\规定分布与向量尺度};
 \node[flow,text width=42mm,right=of coset] (sample) {短原像样本$\mathbf x$\\满足目标分布要求};
 \node[check,text width=42mm,below=of sample] (checks) {模方程、维数与范数检查\\验证必要输出条件};
 \draw[flowarrow](relation)--(coset);
 \draw[flowarrow](coset)--(sample);
 \draw[flowarrow](trapdoor)--(sample);
 \draw[flowarrow](parameter)--(sample);
 \draw[flowarrow](sample)--(checks);
\end{tikzpicture}
\caption{目标陪集、陷门、Gaussian参数与短原像之间的关系}
\label{fig:preimage-sampling-foundation}
\end{figure}

Sageloli、Lai和Wang--PW的采样任务均包含公开目标关系、秘密辅助信息、目标分布与输出检查，具体陷门和采样算法分别在第四、五章说明。

\section{频域表示与低秩线性代数}

\subsection{Gram矩阵与低维子空间}

对复矩阵$J\in\mathbb C^{r\times k}$，$J^*$表示共轭转置，
\[
 G_J=J^*J\in\mathbb C^{k\times k}
\]
称为$J$的Gram矩阵。对任意$\mathbf y$都有
$\mathbf y^*G_J\mathbf y=\lVert J\mathbf y\rVert_2^2\geq0$，所以$G_J$为半正定矩阵。它记录
$J$各列之间的内积关系；$J^*J$与$JJ^*$的非零特征值相同。当$k\ll r$时，分析
$k\times k$的$J^*J$通常比展开$r\times r$的$JJ^*$更经济。

若$J^*J=U\operatorname{diag}(\lambda_i)U^*$，特征向量给出$J$列空间中的正交方向，特征值描述
这些方向的伸缩程度。与$\operatorname{span}(J)$正交的向量$\mathbf x$满足$J^*\mathbf x=0$，
因此任何形如$JJ^*$的修正项都不会作用在该正交补上。这是后文把大空间运算转入低维子空间的
基本依据。

\subsection{Gaussian参数矩阵与平方根作用}

对正定矩阵$\Sigma$，非球形Gaussian可以写成
\[
 \rho_{\sqrt\Sigma}(\mathbf x)
 =\exp(-\pi\mathbf x^*\Sigma^{-1}\mathbf x).
\]
其中$\Sigma$为Gaussian平方参数矩阵；在上述$\pi$归一化下，连续概率分布的协方差为$\Sigma/(2\pi)$。

若$L L^*=\Sigma$，把标准连续Gaussian向量经过$L$变换即可得到由$\Sigma$描述的相关方向。
当
\[
 \Sigma=aI-cJJ^*,\qquad a,c>0,
\]
修正只发生在$J$的列空间；其正交补仍按$\sqrt a$缩放，而列空间内各方向根据
$a-c\lambda_i$调整。因而可以通过较小的Gram矩阵$J^*J$构造平方根作用，而无须保存完整
$r\times r$矩阵。离散陪集采样依次完成连续相关扰动、离散化、目标陪集采样和精确方程检查。第五章的Sageloli结构化实现正是把这一低秩关系逐频率应用到环矩阵。

\section{身份基加密与分层身份基加密}

\subsection{身份基加密}

身份基加密（Identity-Based Encryption，IBE）使用身份字符串代替传统用户公钥。可信发钥
机构运行$\mathsf{Setup}$得到主公钥$\mathsf{mpk}$和主私钥$\mathsf{msk}$；用户通过
$\mathsf{Extract}(\mathsf{msk},\mathsf{id})$领取身份私钥；发送者以$\mathsf{id}$和
$\mathsf{mpk}$加密；接收者使用对应身份私钥解密。其正确性要求对任意合法身份和消息，诚实
生成的身份私钥能够恢复消息。

IBE把公钥发现问题转换为身份规范与发钥机构管理问题。通信者使用对方规范身份和系统主公钥进行加密，发钥机构负责核验身份、保护主私钥并管理密钥生命周期。因而，IBE适合具有明确组织
边界和统一身份管理的系统。

\subsection{HIBE的分层身份与密钥委派}

分层身份基加密（Hierarchical Identity-Based Encryption，HIBE）把身份扩展为路径
$(\mathsf{id}_1,\ldots,\mathsf{id}_d)$。除$\mathsf{Setup}$、$\mathsf{Extract}$、
$\mathsf{Encrypt}$和$\mathsf{Decrypt}$外，HIBE还包含密钥委派$\mathsf{Derive}$：持有前缀
路径密钥的实体可以为下一层身份生成子密钥，无需再次调用根主私钥。

本项目使用两层HIBE。第一层是接收身份$r$，第二层是发送身份$s$：
\[
  (\mathsf{id}_1,\mathsf{id}_2)=(r,s).
\]
发钥机构为$r$生成一级密钥$\mathsf{dk}_r$，接收者再根据认可的发送身份$s$委派叶密钥
$\mathsf{dk}_{r,s}$。发送者以同一身份路径加密。路径中任一身份改变，生成的叶密钥都不能按
正确路径恢复密文。

除消息机密性外，IB-ME还需要HIBE具备身份匿名性。匿名性要求攻击者不能仅从密文判断其对应的
身份路径。该性质用于避免公开密文直接泄露“谁向谁发送”的关系。各HIBE方案的安全模型与参数见第四章。

\section{身份基签名}

一套身份基签名包含四个核心算法：
\[
\begin{aligned}
 (\mathsf{ipk},\mathsf{imsk})&\leftarrow\mathsf{Setup}(1^\lambda),\\
 \mathsf{sk}_{\mathsf{id}}&\leftarrow\mathsf{KeyExtract}(\mathsf{imsk},\mathsf{id}),\\
 \sigma&\leftarrow\mathsf{Sign}(\mathsf{sk}_{\mathsf{id}},M),\\
 b&\leftarrow\mathsf{Verify}(\mathsf{ipk},\mathsf{id},M,\sigma).
\end{aligned}
\]
正确性要求合法身份密钥产生的签名能够通过验证。不可伪造性要求攻击者即使获得其他身份的密钥
和多条签名，也不能为未获授权的身份与新消息生成有效签名。

IBS可以采用原生身份关系，也可以由普通数字签名经认证转换得到。认证转换中，发钥机构为用户
普通验签公钥签发与身份绑定的证书，完整IBS签名携带用户公钥、机构证书和用户签名；验证者先
检查证书，再检查用户签名。原生格IBS则让身份直接进入目标向量、挑战或陷门采样关系。两类方案
在接口上都完成“身份＋消息＋签名”的公开验证，但密钥尺寸、签名尺寸和运行成本不同。

在IB-ME中，IBS的职责是认证发送者。接收端使用自己认可的发送身份验签，因而更换$s$不仅会
改变HIBE第二层路径，也会改变IBS验证输入。两种检查共同把发送身份绑定到消息恢复过程。

\section{身份基匹配加密与通用构造}

IB-ME包含系统建立、发送者发钥、接收者发钥、消息加密和匹配解密五项算法。发送者发钥为身份
$s$生成加密所需的发送方密钥，接收者发钥为身份$r$生成接收方密钥。加密算法同时接收发送方密钥、
接收身份$r$和消息$M$；解密算法同时接收接收方密钥、认可发送身份$s'$和密文。正确性条件为
$r=r'$且$s=s'$时恢复$M$。

Wang等人的通用构造使用两层匿名HIBE与IBS实现这一语法\cite{wang2024ibme}。HIBE沿$(r,s)$
路径保护内部秘密，IBS以$s$认证发送者生成的协议消息。两种原语承担不同职责：HIBE决定哪一条
双方身份路径能够恢复秘密，IBS决定恢复后的消息是否确由指定发送身份产生。通用构造不要求二者
使用相同的维度、模数或格结构。

该构造讨论的基本目标包括消息与身份隐私、发送来源认证以及合法用户之间的自我伪造防护。软件通过完整密文字节认证、规范解析和状态恢复组织消息处理，五个公共算法经原语接口调用HIBE与IBS。

\subsection{安全目标与证明模型术语}

正确性说明诚实生成的密钥和密文能否按规定恢复消息；安全性则要在明确的攻击者能力下描述其
不能完成的任务。与本作品相关的常见目标包括：加密方案的消息机密性与身份匿名性，签名方案的
选择消息攻击下存在性不可伪造（EUF-CMA），以及把新签名值也纳入限制的强不可伪造
（sEUF-CMA）。IB-ME还要把接收身份、发送身份和合法用户持钥能力纳入隐私与认证实验。第四章按方案列出安全模型及其原论文依据
\cite{wang2024ibme,sageloli2023ibs}。

随机预言机模型（Random Oracle Model，ROM）把哈希函数抽象为可查询的理想随机函数；量子
随机预言机模型（Quantum Random Oracle Model，QROM）进一步允许攻击者以量子叠加方式查询
该预言机。标准模型构造不依赖这种理想随机预言机。ROM、QROM和标准模型分别规定安全证明中的哈希抽象与攻击者查询能力。

\section{本章小结}

本章从格、模运算和短向量出发，介绍了LWE、SIS及其环、模和NTRU结构变体，区分$q$元核格与
目标陪集，并说明离散Gaussian、gadget矩阵和结构化陷门如何共同支持短原像采样。环伴随、频域
表示、Gram矩阵与低秩平方根进一步解释了结构化计算为何能够避免稠密展开。最后给出HIBE、IBS
和IB-ME的基本语法，以及正确性、安全目标和证明模型之间的区别。
